%%% FIELD prop-stationary-shared-response-endpoints-proof
Write \(e=(e_0,\ldots,e_T)\), and define the finite set of factual paths compatible with the coarsened evidence by
\[
 \mathcal H_e
 =\{x_{0:T}\in\mathcal S^{T+1}:o(x_t)=e_t\text{ for every }t=0,\ldots,T\}.
\]
For \(q\in\Theta(P)\), \(x_{0:T}\in\mathcal H_e\), \(z_{0:T}\in\mathcal S^{T+1}\), and \(t\geq 1\), put
\[
 A_t^q(x_{0:T},z_{0:T})
 =\sum_{f\in\mathcal R}q_f
 \mathbf 1\!\left\{f(x_{t-1},b(x_{t-1}))=x_t,\ 
 f(z_{t-1},\pi(z_{t-1}))=z_t\right\}.
\]
The shared-origin condition in \(\mathfrak M_{d,k,T}(\epsilon,B)\), the coupled-trajectory construction, initial exogeneity, and the iid-response condition imply, date by date, that
\[
 \Pr_q(X_{0:T}=x_{0:T},Z_{0:T}=z_{0:T})
 =\rho(x_0)\mathbf 1\{z_0=x_0\}
   \prod_{t=1}^T A_t^q(x_{0:T},z_{0:T}).                                      \tag{1}
\]
Indeed, after the common initial state is fixed, the two transitions at date \(t\) are deterministic functions of the same draw \(F_t\), so their joint conditional probability is exactly \(A_t^q\); independence of the draws across dates gives the product in (1).

Define the finite path-sum numerator
\[
 N_e(q)=
 \sum_{x_{0:T}\in\mathcal H_e}
 \sum_{z_{0:T}\in\mathcal S^{T+1}}
 \rho(x_0)\mathbf 1\{z_0=x_0\}
 \left(\frac1T\sum_{s=1}^T r(z_s)\right)
 \prod_{t=1}^T A_t^q(x_{0:T},z_{0:T}).                                      \tag{2}
\]
This is precisely the joint expectation of \(T^{-1}\sum_{s=1}^T r(Z_s)\) times \(\mathbf 1\{E_T=e\}\), by (1) and the definition of \(\mathcal H_e\).

It remains to identify the denominator.  Marginally in the factual world, initial exogeneity and the iid-response condition give
\[
 \Pr_q(X_{0:T}=x_{0:T})
 =\rho(x_0)\prod_{t=1}^T
   \sum_{f\in\mathcal R}q_f
   \mathbf 1\{f(x_{t-1},b(x_{t-1}))=x_t\}.                                 \tag{3}
\]
Because \(q\in\Theta(P)\), each factor in (3) equals
\[
 \sum_{f\in\mathcal R}q_f
 \mathbf 1\{f(x_{t-1},b(x_{t-1}))=x_t\}
 =P(x_t\mid x_{t-1},b(x_{t-1})).                                           \tag{4}
\]
Summing (3) over \(x_{0:T}\in\mathcal H_e\) and using (4) and the definition of the factual-record law therefore yields
\[
 \Pr_q(E_T=e)
 =\sum_{x_{0:T}\in\mathcal H_e}\rho(x_0)
   \prod_{t=1}^T P(x_t\mid x_{t-1},b(x_{t-1}))
 =\Pr_{P,\rho,b,o}(E_T=e)=:p_e.                                            \tag{5}
\]
Thus the evidence probability is the same for every compatible mechanism.  The stated evidence-support condition gives \(p_e>0\), which justifies division, and (2)--(5) prove the explicit identity
\[
 Q_{q,T}(P,\rho,b,o,e)=\frac{N_e(q)}{p_e}.                                  \tag{6}
\]
Equations (2) and (6) are exactly the objective evaluated by the stationary shared-response endpoint program: every factual path compatible with \(e\) and every alternative path is enumerated, the displayed joint-response factor is multiplied at every date, and the resulting reward numerator is divided by the common factual evidence probability.

We next prove that the displayed minimum and maximum exist.  The response-program set \(\mathcal R=\mathcal S^{\mathcal S\times\mathcal A}\) is finite.  The model-class membership of the tuple supplies the displayed \(q\in\Theta(P)\), so \(\Theta(P)\) is nonempty.  It is the subset of the finite-dimensional simplex specified by the finitely many linear equalities in its definition, and hence is closed and bounded.  More explicitly, every sequence in \(\Theta(P)\) has, by repeated subsequence extraction over the finitely many coordinates in \([0,1]\), a coordinatewise convergent subsequence; the limiting vector remains nonnegative, has total mass one, and satisfies all compatibility equalities after taking limits.  Thus its limit lies in \(\Theta(P)\).

Each \(A_t^q(x_{0:T},z_{0:T})\) is linear in \(q\), so (2) is a polynomial of degree at most \(T\) in the finitely many coordinates \((q_f)_{f\in\mathcal R}\).  By \(p_e>0\), (6) is continuous on \(\Theta(P)\).  Let
\[
 m=\inf_{q\in\Theta(P)}Q_{q,T}(P,\rho,b,o,e),\qquad
 M=\sup_{q\in\Theta(P)}Q_{q,T}(P,\rho,b,o,e).
\]
The objective lies in \([0,1]\) because \(r\) does.  Choose sequences whose objective values tend to \(m\) and \(M\).  The subsequence property just proved and continuity of (6) provide limits \(q^-,q^+\in\Theta(P)\) satisfying
\[
 Q_{q^-,T}(P,\rho,b,o,e)=m,
 \qquad
 Q_{q^+,T}(P,\rho,b,o,e)=M.                                                \tag{7}
\]

Finally, \(\Theta(P)\) is convex because all its defining constraints are linear.  If \(m<M\), then for \(s\in[0,1]\),
\[
 q_s=(1-s)q^-+s q^+\in\Theta(P),
 \qquad
 \varphi(s)=Q_{q_s,T}(P,\rho,b,o,e)
\]
is a continuous polynomial quotient with \(\varphi(0)=m\) and \(\varphi(1)=M\).  For any \(v\in(m,M)\), the elementary nested-interval construction that retains endpoints \(a_n,b_n\) with \(\varphi(a_n)\leq v\leq\varphi(b_n)\) and bisects \([a_n,b_n]\) produces a common limit \(s_v\); continuity gives \(\varphi(s_v)=v\).  If \(m=M\), the assertion is immediate.  Consequently,
\[
 \{Q_{q,T}(P,\rho,b,o,e):q\in\Theta(P)\}=[m,M].                            \tag{8}
\]
By the definition of the sharp interval, (7)--(8) show that the endpoint program returns exactly \(I_T(P,\rho,b,o,e)\), with attained endpoints.  In (1)--(2) the optimizer \(q^-\), respectively \(q^+\), is one fixed response-mass vector appearing in every factor, at every date, and for every hidden path; no date-specific or branch-specific selector has been introduced.  This proves the final shared-stationarity assertion.  The positive-cell and Poisson conditions are part of the stated model class but are not needed for this finite-horizon endpoint characterization beyond the separately assumed positivity of the realized evidence probability.
